

%------------------------------------------------------------
\begin{frame}{Indexation de données géométriques}

Certaines applications stockent des données géométriques
(ex, cadastre, réseaux routiers, etc.). Requêtes
courantes\,:
\begin{enumerate}
  \item pointé,
   \item fenêtrage,
  \item jointure spatiale. 
\end{enumerate}
Comment optimiser ces opérations avec des index\,?

\end{frame}
%------------------------------------------------------------
\begin{frame}{Obstacles}

Principaux obstacles à l'utilisation d'un index standard\,:
\begin{enumerate}
  \item On ne peut pas utiliser un arbre-B

      car 
      pas d'ordre sur les points ou les polygones.

  \item Opérations coûteuses

      Algorithmes géométriques, souvent de complexité
            polynomiale, difficiles à implanter.
\end{enumerate}

$\Rightarrow$ des structures spécialisées pour les
données géométriques

\end{frame}
%------------------------------------------------------------
\begin{frame}{Principes}

Pour limiter le coût des opérations, on 
utilise une approximation rectangulaire 
des données géométriques\,: le rectangle englobant.

\begin{enumerate}
  \item l'index est construit sur l'ensemble des rectangles\,;
  \item On effectue un {\em filtrage} basé sur les rectangles
         pour les pointés, fenêtrages et jointures.
  \item Quand on accède aux objets eux-mêmes, on effectue
          une seconde opération sur la vraie géométrie
          (étape de {\em raffinement})
\end{enumerate}


\end{frame}
%------------------------------------------------------------
\begin{frame}{L'arbre R, définition}

Assez proche de l'arbre-B dans ses principes\,:

\begin{enumerate}
  \item c'est un arbre équilibré\,;
  \item chaque entrée est une paire $[adresse, rectangle]$\,;
  \item tous les n{\oe}uds sont occupés au moins à 50\,\%
  \item le rectangle d'une entrée  $[adresse, rectangle]$
          englobe {\em tous} les rectangles du n{\oe}ud
           pointé par $adresse$\,;
\end{enumerate}


\end{frame}
%------------------------------------------------------------
\begin{frame}{L'arbre R, exemple}


\figSlideWithSize{rtree}{11cm}


\end{frame}
%------------------------------------------------------------
\begin{frame}{L'arbre R, pointé}


\figSlideWithSize{pqRtree}{11cm}


\end{frame}
%------------------------------------------------------------
\begin{frame}{L'arbre R, insertion (1)}


\figSlideWithSize{insRtree1}{11cm}


\end{frame}
%------------------------------------------------------------
\begin{frame}{L'arbre R, insertion (2)}


\figSlideWithSize{insRtree2}{11cm}


\end{frame}
